\documentclass[11pt,a4paper]{article}

% --- Idioma y codificación ---
\usepackage[utf8]{inputenc}
\usepackage[T1]{fontenc}
\usepackage[spanish]{babel}

% --- Formato de página ---
\usepackage[a4paper,margin=2.5cm]{geometry}
\usepackage{setspace}
\setstretch{1.05}

% --- Tablas y columnas ---
\usepackage{array}
\usepackage{tabularx}
\usepackage{longtable}
\usepackage{booktabs}

% --- Listas ---
\usepackage{enumitem}

% --- Encabezados y pies ---
\usepackage{fancyhdr}
\usepackage{lastpage}

% --- Símbolos (checkbox) ---
\usepackage{pifont}
\newcommand{\checkbox}{\ding{113}} % cuadro vacío
\newcommand{\checked}{\ding{51}}   % check

% --- Enlaces (email) ---
\usepackage[hidelinks]{hyperref}

% --- Ajustes ---
\setlength{\parindent}{0pt}
\setlength{\parskip}{6pt}

% Encabezado/pie para páginas 2+
\pagestyle{fancy}
\fancyhf{}
\fancyhead[L]{Geoinformática - USACH}
\fancyhead[C]{Primera Evaluación de Proyecto}
\fancyhead[R]{2025-1}
\fancyfoot[L]{Prof. Francisco Parra O.}
\fancyfoot[C]{Página \thepage\ de 10}
\fancyfoot[R]{23 de septiembre de 2025}
\renewcommand{\headrulewidth}{0.4pt}
\renewcommand{\footrulewidth}{0.4pt}

\begin{document}

% --- Primera página (sin encabezado) ---
\thispagestyle{plain}

\begin{center}
\textbf{UNIVERSIDAD DE SANTIAGO DE CHILE}\\
Facultad de Ingeniería\\
Departamento de Ingeniería Informática\\[10pt]

\textbf{\Large PRIMERA EVALUACION DE PROYECTO}\\[6pt]
\textbf{\large Análisis Geoespacial Aplicado}\\
\textbf{\large Presentación de Avances e Informe Técnico}\\[10pt]

Prof. Francisco Parra O.\\
\href{mailto:francisco.parra@usach.cl}{francisco.parra@usach.cl}\\[8pt]

Fecha de publicación: 23 de septiembre de 2025\\
Fecha de presentaciones: 7 de octubre de 2025\\
Fecha de entrega del informe: 10 de octubre de 2025
\end{center}

\section{Descripción General}
Esta primera evaluación constituye un hito fundamental en el desarrollo del proyecto semestral de su grupo. Después de 4 semanas desde la selección del proyecto (2 ya transcurridas + 2 adicionales), deben demostrar avances concretos en:

\begin{itemize}[leftmargin=1.2cm]
  \item Comprensión del problema: Marco teórico y justificación clara
  \item Adquisición de datos: Fuentes identificadas y datos descargados
  \item Análisis exploratorio: Primeras visualizaciones y estadísticas
  \item Pipeline técnico: Ambiente de desarrollo funcional
  \item Plan de trabajo: Cronograma realista para el resto del semestre
\end{itemize}

\section{Objetivos de la Evaluación}

\subsection{Objetivos Principales}
\begin{enumerate}[leftmargin=1.2cm]
  \item Verificar factibilidad: Confirmar que el proyecto es realizable con los datos y tiempo disponibles
  \item Evaluar comprensión: Asegurar que entienden el problema geoespacial a resolver
  \item Validar enfoque técnico: Revisar que las herramientas y métodos sean apropiados
  \item Proporcionar retroalimentación: Guiar ajustes necesarios antes de continuar
\end{enumerate}

\newpage

\section{Componentes de la Evaluación}

\subsection{Distribución de Ponderaciones}
\begin{table}[h!]
\centering
\begin{tabularx}{\linewidth}{@{}l c X@{}}
\toprule
\textbf{Componente} & \textbf{Peso} & \textbf{Descripción}\\
\midrule
Informe Técnico & 40 \% & Documento escrito con todos los requisitos\\
Presentación Oral & 25 \% & Exposición de 15 minutos + 5 de preguntas\\
Código y Reproducibilidad & 20 \% & Repositorio GitHub con código funcional\\
Visualizaciones & 10 \% & Calidad de mapas y gráficos producidos\\
Trabajo en Equipo & 5 \% & Evidencia de colaboración efectiva\\
\midrule
\textbf{Total} & \textbf{100 \%} & \\
\bottomrule
\end{tabularx}
\caption{Ponderación de componentes de evaluación}
\end{table}

\section{Requisitos Mínimos por Componente}

\subsection{A. Informe Técnico (40 \%)}
\textbf{Requisitos Mínimos}\\
Extensión: 15--20 páginas (sin incluir anexos)\\
Formato: LaTeX o Markdown convertido a PDF\\
Estructura obligatoria:

\subsubsection{Estructura del Informe}
\begin{enumerate}[leftmargin=1.2cm]
  \item \textbf{Resumen Ejecutivo (1 página)}
  \begin{itemize}[leftmargin=1.2cm]
    \item Problema abordado
    \item Área de estudio
    \item Metodología propuesta
    \item Resultados preliminares
  \end{itemize}

  \item \textbf{Introducción (2--3 páginas)}
  \begin{itemize}[leftmargin=1.2cm]
    \item Contexto del problema
    \item Justificación de la relevancia
    \item Objetivos específicos
    \item Preguntas de investigación
  \end{itemize}

  \item \textbf{Marco Teórico (3--4 páginas)}
  \begin{itemize}[leftmargin=1.2cm]
    \item Revisión de literatura (mínimo 10 referencias)
    \item Conceptos geoespaciales clave
    \item Casos de estudio similares
    \item Estado del arte en la temática
  \end{itemize}

  \item \textbf{Área de Estudio (2--3 páginas)}
  \begin{itemize}[leftmargin=1.2cm]
    \item Delimitación geográfica precisa
    \item Características territoriales relevantes
    \item Mapa de ubicación (obligatorio)
    \item Justificación de la selección
  \end{itemize}

  \item \textbf{Datos y Metodología (4--5 páginas)}
  \begin{itemize}[leftmargin=1.2cm]
    \item Fuentes de datos:
    \begin{itemize}[leftmargin=1.2cm]
      \item Tabla con todas las fuentes
      \item Descripción de cada dataset
      \item Fechas de actualización
      \item Resolución espacial/temporal
    \end{itemize}
    \item Procesamiento preliminar:
    \begin{itemize}[leftmargin=1.2cm]
      \item Diagrama de flujo metodológico
      \item Herramientas utilizadas
      \item Transformaciones aplicadas
    \end{itemize}
  \end{itemize}

  \item \textbf{Resultados Preliminares (3--4 páginas)}
  \begin{itemize}[leftmargin=1.2cm]
    \item Análisis exploratorio de datos (EDA)
    \item Estadísticas descriptivas espaciales
    \item Primeras visualizaciones
    \item Patrones identificados
  \end{itemize}

  \item \textbf{Cronograma (1 página)}
  \begin{itemize}[leftmargin=1.2cm]
    \item Gantt chart con actividades restantes
    \item Hitos principales
    \item Distribución de responsabilidades
  \end{itemize}

  \item \textbf{Conclusiones Preliminares (1 página)}
  \begin{itemize}[leftmargin=1.2cm]
    \item Factibilidad del proyecto
    \item Desafíos identificados
    \item Ajustes propuestos
  \end{itemize}
\end{enumerate}

\subsection{B. Presentación Oral (25 \%)}
\textbf{Entregables}\\
Duración: 15 minutos de exposición + 5 minutos de preguntas\\
Formato: Presencial con apoyo de slides\\
Participación: Todos los integrantes del grupo deben participar en la exposición

\subsubsection{Estructura de la Presentación}
\begin{enumerate}[leftmargin=1.2cm]
  \item \textbf{Introducción (2 min)}
  \begin{itemize}[leftmargin=1.2cm]
    \item Problema y motivación
    \item Objetivos del proyecto
  \end{itemize}

  \item \textbf{Área de Estudio (2 min)}
  \begin{itemize}[leftmargin=1.2cm]
    \item Mapas de localización
    \item Características relevantes
  \end{itemize}

  \item \textbf{Datos y Métodos (4 min)}
  \begin{itemize}[leftmargin=1.2cm]
    \item Fuentes de datos (con ejemplos visuales)
    \item Pipeline de procesamiento
    \item Herramientas utilizadas
  \end{itemize}

  \item \textbf{Resultados Preliminares (5 min)}
  \begin{itemize}[leftmargin=1.2cm]
    \item Demo en vivo (opcional pero valorado)
    \item Visualizaciones principales
    \item Hallazgos iniciales
  \end{itemize}

  \item \textbf{Próximos Pasos (2 min)}
  \begin{itemize}[leftmargin=1.2cm]
    \item Plan de trabajo
    \item Desafíos esperados
  \end{itemize}
\end{enumerate}

\subsection{C. Código y Reproducibilidad (20 \%)}
\textbf{Requisitos Mínimos}\\
Repositorio GitHub: \texttt{github.com/[usuario]/geoinformatica-proyecto}\\
Estructura mínima requerida:

\begin{verbatim}
proyecto/
README.md # Documentación principal
requirements.txt # Dependencias Python
data/
  raw/ # Datos originales (o links de descarga)
  processed/ # Datos procesados
notebooks/
  01_exploratory.ipynb # Análisis exploratorio
  02_preprocessing.ipynb
  03_analysis.ipynb
src/
  data_download.py # Scripts de descarga
  preprocessing.py # Funciones de procesamiento
  visualization.py # Funciones de visualización
outputs/
  figures/ # Gráficos generados
  maps/ # Mapas producidos
docs/
  informe_v1.pdf # Informe de primera entrega
\end{verbatim}

\subsubsection{Requisitos de Código}
\begin{itemize}[leftmargin=1.2cm]
  \item Documentación: Cada función debe tener docstring
  \item Reproducibilidad: Instrucciones claras para ejecutar
  \item Datos: Si son $\ge 100$MB, incluir script de descarga
  \item Ambiente: Archivo \texttt{requirements.txt} o \texttt{environment.yml}
\end{itemize}

\subsection{D. Visualizaciones (10 \%)}
\textbf{Criterios de Evaluación}\\
Se evaluará la calidad técnica y comunicativa de las visualizaciones

\subsubsection{Requisitos Mínimos de Visualización}
\textbf{3 mapas temáticos como mínimo:}
\begin{itemize}[leftmargin=1.2cm]
  \item Mapa de ubicación del área de estudio
  \item Mapa de datos principales
  \item Mapa de análisis o resultado preliminar
\end{itemize}

\textbf{5 gráficos estadísticos:}
\begin{itemize}[leftmargin=1.2cm]
  \item Histogramas de variables clave
  \item Series temporales (si aplica)
  \item Correlaciones espaciales
  \item Diagramas de dispersión
\end{itemize}

\textbf{1 visualización interactiva} (Folium, Plotly, o Streamlit)

\section{Checklist de Entregables}
\textbf{Importante}\\
Todos los items de esta lista son obligatorios. La ausencia de cualquiera resultará en descuento de puntaje.

\subsection{Checklist Completo}
\begin{table}[h!]
\centering
\begin{tabularx}{\linewidth}{@{}c X l@{}}
\toprule
 & \textbf{Item} & \textbf{Estado}\\
\midrule
\checkbox & Informe en formato PDF (15--20 páginas) & Pendiente\\
\checkbox & Repositorio GitHub público y documentado & Pendiente\\
\checkbox & Mínimo 3 mapas temáticos con elementos cartográficos & Pendiente\\
\checkbox & Mínimo 5 gráficos estadísticos & Pendiente\\
\checkbox & 1 visualización interactiva funcional & Pendiente\\
\checkbox & Datos descargados y organizados & Pendiente\\
\checkbox & Código Python ejecutable y documentado & Pendiente\\
\checkbox & Análisis exploratorio completo (EDA) & Pendiente\\
\checkbox & Marco teórico con mínimo 10 referencias & Pendiente\\
\checkbox & Cronograma detallado (Gantt) & Pendiente\\
\checkbox & Presentación slides (PDF o PPTX) & Pendiente\\
\checkbox & README.md con instrucciones de reproducción & Pendiente\\
\checkbox & requirements.txt o environment.yml & Pendiente\\
\checkbox & Evidencia de trabajo colaborativo (commits de todos los integrantes) & Pendiente\\
\checkbox & Diagrama de flujo metodológico & Pendiente\\
\bottomrule
\end{tabularx}
\end{table}

\section{Rúbrica de Evaluación Detallada}

\subsection{Niveles de Desempeño}
\begin{table}[h!]
\centering
\small
\begin{tabularx}{\linewidth}{@{}l X X X X@{}}
\toprule
\textbf{Criterio} &
\textbf{Excelente (7.0)} &
\textbf{Bueno (6.0)} &
\textbf{Suficiente (5.0)} &
\textbf{Insuficiente ($<4.0$)}\\
\midrule
Comprensión del Problema &
Demuestra comprensión profunda, identifica complejidades y propone soluciones innovadoras &
Comprende bien el problema y propone soluciones adecuadas &
Comprensión básica con soluciones estándar &
Comprensión limitada o incorrecta\\
\addlinespace
Calidad de Datos &
Múltiples fuentes integradas, datos actualizados, metadata completa &
Buenas fuentes, datos apropiados &
Datos mínimos necesarios &
Datos insuficientes o inadecuados\\
\addlinespace
Análisis Técnico &
Análisis sofisticado, métodos avanzados bien aplicados &
Buen análisis con métodos apropiados &
Análisis básico pero correcto &
Análisis superficial o con errores\\
\addlinespace
Visualizaciones &
Excepcionales, claras, estéticamente superiores, interactivas &
Buena calidad, claras y bien diseñadas &
Correctas pero básicas &
Pobres o con errores\\
\addlinespace
Código &
Limpio, bien documentado, modular, eficiente &
Bien estructurado y funcional &
Funciona pero poco organizado &
Desorganizado o no funciona\\
\addlinespace
Presentación &
Clara, fluida, dominio del tema, responde preguntas con solidez &
Buena presentación, responde bien &
Presentación correcta &
Confusa o incompleta\\
\bottomrule
\end{tabularx}
\caption{Rúbrica de evaluación por criterios}
\end{table}

\section{Ejemplos de Proyectos Exitosos}

\subsection{Características de Proyectos Destacados}
Para alcanzar la excelencia (nota 6.5--7.0), los proyectos anteriores han incluido:

\begin{itemize}[leftmargin=1.2cm]
  \item Integración de múltiples fuentes: OSM + Census + Sentinel + DEM
  \item Análisis espacial avanzado: Hot spots, clusters, interpolación
  \item Machine Learning básico: Clasificación o predicción espacial
  \item Dashboard interactivo: Streamlit con mapas Folium
  \item Reproducibilidad total: Docker o ambiente virtual completo
\end{itemize}

Impacto social: Problema real con aplicación práctica

\section{Recursos de Apoyo}

\subsection{Herramientas Recomendadas}
\textbf{Adquisición de datos:}
\begin{itemize}[leftmargin=1.2cm]
  \item OSMnx para OpenStreetMap
  \item Google Earth Engine para imágenes satelitales
  \item IDE Chile para datos oficiales
  \item Census API para datos demográficos
\end{itemize}

\textbf{Procesamiento:}
\begin{itemize}[leftmargin=1.2cm]
  \item GeoPandas para vectores
  \item Rasterio para rasters
  \item PySAL para análisis espacial
\end{itemize}

\textbf{Visualización:}
\begin{itemize}[leftmargin=1.2cm]
  \item Folium para mapas web
  \item Matplotlib/Seaborn para gráficos
  \item Plotly para gráficos interactivos
  \item Streamlit para dashboards
\end{itemize}

\subsection{Plantillas Disponibles}
En el repositorio del curso encontrarán:
\begin{itemize}[leftmargin=1.2cm]
  \item Template LaTeX para el informe
  \item Estructura base de proyecto Python
  \item Ejemplos de notebooks bien documentados
  \item Scripts de descarga de datos
\end{itemize}

\section{Fechas Importantes}
\begin{itemize}[leftmargin=1.2cm]
  \item Publicación de esta guía: 23 de septiembre de 2025
  \item Consultas y dudas: Durante las próximas 2 semanas en clase
  \item Presentaciones: 7 de octubre de 2025
  \item Entrega de informe y código: 10 de octubre de 2025
  \item Retroalimentación: Una semana después de las presentaciones
\end{itemize}

\section{Criterios de Penalización}
\textbf{Importante}\\
Las siguientes situaciones resultarán en penalizaciones:
\begin{itemize}[leftmargin=1.2cm]
  \item Entrega tardía: -1.0 punto por día de atraso
  \item Falta de algún integrante a la presentación: -0.5 puntos
  \item Código no reproducible: -0.5 puntos
  \item Sin commits de algún integrante: -0.3 puntos
  \item Plagio o copia: Nota 1.0 y revisión por comité de ética
\end{itemize}

\section{Preguntas Frecuentes}

\subsection{¿Qué pasa si nuestros datos aún no están completos?}
Deben mostrar avance con los datos disponibles y explicar claramente el plan para obtener los faltantes. Es mejor ser transparente sobre las limitaciones.

\subsection{¿Podemos cambiar el enfoque del proyecto?}
Sí, pero debe justificarse en el informe. Esta evaluación es precisamente para detectar si se necesitan ajustes.

\subsection{¿Qué nivel de análisis espacial se espera?}
Como mínimo: estadísticas descriptivas espaciales, mapas de densidad/distribución, y algún índice espacial (Moran's I, Getis-Ord, etc.).

\subsection{¿El código debe estar 100 \% terminado?}
No, pero debe ser funcional para los análisis presentados y estar bien documentado. Se espera un 40--50 \% de avance del proyecto total.

\section{Contacto y Consultas}
Horarios de consulta:\\
Presencial: Horario a coordinar\\
Zoom: Previa coordinación por email\\
Email: \href{mailto:francisco.parra@usach.cl}{francisco.parra@usach.cl}

\vfill
\begin{center}
\textbf{¡Éxito en su primera evaluación!}\\
Prof. Francisco Parra O.
\end{center}

\end{document}
